The unreliability of reward signals is one of the core challenges for deep reinforcement learning in real-world applications. In settings such as sim-to-real fine-tuning and long-horizon robotic manipulation, reward signals are either extremely sparse, provided only upon task completion, or not fully aligned with the true objective, and therefore cannot reliably steer behavior toward the correct outcome. This thesis systematically compares, on the Door task of HumanoidBench, the fine-tuning behaviour of PPO, which uses GAE-style absolute advantage estimation, and GRPO, which uses group-relative advantage estimation, under two reward conditions: a Misaligned Reward condition that simulates unreliable reward information, and a Sparse Reward condition that simulates extremely sparse rewards.

Our experimental results show that under the Misaligned condition PPO's return soars from a baseline of 244 to 431 while the success rate drops from 5.0\% to 0.0\%, exhibiting systematic reward hacking; under the Sparse condition PPO either collapses or stalls. In contrast, GRPO remains stable under both conditions, with the return held within a narrow band of 188--219. KL divergence and parameter-space distance analyses show that GRPO's policy update is not stagnant, yet its return is decoupled from KL: the update magnitude varies with seed and condition while the return stays stable, whereas PPO's KL and return are tightly coupled. This indicates that GRPO's update magnitude is insensitive to the scale and density of the reward signal, and that its group-relative advantage mechanism may inherently suppress the positive-feedback loop of reward hacking, making GRPO a safer fine-tuning strategy when reward design is imperfect or rewards are extremely sparse.
